\subsection{Hurwitz 移动}\label{sec:18i}
设 $\gamma_k,\gamma_{k+1},\delta_k$ 为\sr{fig:16.2}{图~16.2} 所示的消失路径，$\widetilde\Delta_k^\#,\widetilde\Delta_{k+1}^\#,\widetilde\Pi_k^\#$ 为相应的 (B) 型膜。更准确地说，$\widetilde\Delta_k^\#$ 的分次任意，但其余两个分次应作调整，使它们与 Dehn 扭转的典范分次一起，与\stref{引理}{18.3}中的同痕相容。重述该结果，现在有膜的同痕
\begin{equation}\label{eq:18.28}
\widetilde\Pi_k^\#\simeq\tau_{\widetilde\Delta_k}(\widetilde\Delta_{k+1}^\#).
\end{equation}
如\stref{笔记}{18.5}所述，这一同痕实际经过 $\iota_{\widetilde E}$-不变 Lagrangian 子流形。为使它成为等变膜的同痕，还需比较等变 Pin 结构。由\stref{引理}{14.2}，选择一个等变 Pin 结构，对应于选择不动点集的余定向。我们的约定是，为 $\widetilde\Delta_{k+1}$ 与 $\widetilde\Pi_k$ 选择的余定向，在 $D\widetilde\pi$ 下投影到 $\gamma'_{k+1}(0)$ 和 $\delta'_k(0)$ 的虚部为正的平方根。\znote{286}{这与 §18e 对 (B) 型膜的规定是同一条：等变 Pin 结构对应不动点集的余定向（引理 14.2），而对 $\widetilde\Delta_\gamma$ 有 $D\widetilde\pi(T\widetilde\Delta_\gamma|V_\gamma)=\R\sqrt{\gamma'(0)}$，故「取虚部为正的平方根」唯一确定余定向。要点在于两个球面的余定向分别用 $\gamma'_{k+1}$、$\delta'_k$ 定义，比较之前必须借\eref{18.28}的同痕把二者对齐；笔记 18.5 中相交数 $\widetilde\delta_k\cdot\widetilde\gamma_{k+1}=-1$ 的计算，正是说明这两个选择相容而非相差一个符号。[来源：原书 §18e、§18h、笔记 18.5（印刷页 279、282、286）]}\stref{笔记}{18.5}中的计算说明，这些选择确实与\eref{18.28}相容。现在可应用\secref{17k}的讨论，特别是\stref{引理}{17.23}。将由此得到的代数数据代入\stref{引理}{6.4}，便得到：

\StatementNumber{23}\begin{proposition}\label{st:18.23}
在 $D\F(\pi)$ 中，$\widetilde\Pi_k^\#$ 同构于 $T_{\widetilde\Delta_k^\#}\widetilde\Delta_{k+1}^\#$。
\end{proposition}
